%-------------------------------------------------------------
% LETTRE DE BASE - FRANÇAIS (modèle vierge)
% Configuration (une seule fois) : l'agent remplace chaque [[MOI: ...]]
% avec les infos de l'utilisateur et écrit le paragraphe de réalisations
% à partir de son CV. Il laisse les autres [[...]] : ils se remplissent
% pour chaque offre.
% Pour chaque offre : le skill copie ce fichier dans
% candidatures/<dossier>/lettre.tex, remplace chaque [[...]] et réécrit
% les paragraphes marqués "ADAPTER". FORMAT FIGÉ : ne pas toucher au
% préambule ni à la mise en page. Pas de ligne « Objet ».
% Longueur visée : 250-350 mots, 1 page.
%-------------------------------------------------------------
\documentclass[letterpaper,11pt]{article}

\usepackage[T1]{fontenc}
\usepackage[utf8]{inputenc}
\usepackage{latexsym}
\usepackage[empty]{fullpage}
\usepackage{titlesec}
\usepackage{enumitem}
\usepackage[hidelinks]{hyperref}
\usepackage{fancyhdr}
\usepackage[french]{babel}
\usepackage{tabularx}

\pagestyle{fancy}
\fancyhf{}
\renewcommand{\headrulewidth}{0pt}
\renewcommand{\footrulewidth}{0pt}

\addtolength{\oddsidemargin}{-0.5in}
\addtolength{\evensidemargin}{-0.5in}
\addtolength{\textwidth}{1in}
\addtolength{\topmargin}{-.5in}
\addtolength{\textheight}{1.0in}

\urlstyle{same}

\begin{document}

\begin{flushleft}
    \textbf{\Large [[MOI: Prénom Nom]]} \\
    [[MOI: téléphone]] $|$ \href{mailto:[[MOI: courriel]]}{[[MOI: courriel]]} \\
    [[MOI: Ville, QC]] $|$ \href{[[MOI: URL LinkedIn]]}{\underline{LinkedIn}}
\end{flushleft}

\vspace{20pt}

\begin{flushright}
    [[MOI: Ville]], le \today
\end{flushright}

\vspace{10pt}

\begin{flushleft}
    À l'attention de l'équipe de recrutement \\
    \textbf{[[ENTREPRISE]]} \\
    [[VILLE, PROVINCE]]
\end{flushleft}

\vspace{20pt}

Madame, Monsieur,

\vspace{10pt}

% Intro. Si l'utilisateur fait un stage en ce moment, mentionner « actuellement
% stagiaire chez X » ; s'il est terminé, « ayant effectué un stage chez X ».
Étudiant en [[MOI: programme]] à [[MOI: établissement]], j'ai le plaisir de vous soumettre ma candidature pour le poste de [[TITRE DU POSTE]] au sein de [[ENTREPRISE]] pour [[SESSION, ex. « l'hiver 2027 »]].

\vspace{10pt}

% ADAPTER : 3-4 phrases. Garder les réalisations du CV les plus pertinentes pour
% l'offre et citer 2-4 technos de l'offre que l'utilisateur a réellement utilisées.
% Pas de liste de mots-clés.
[[MOI: Paragraphe par défaut (3-4 phrases) sur les 2-3 meilleures réalisations du CV, avec des résultats concrets.]]

\vspace{10pt}

% ADAPTER : 2-3 phrases sur l'entreprise et le poste, tirées de l'offre
% (mission, produits, responsabilités). Ne rien inventer sur l'entreprise.
[[Ce qui attire l'utilisateur chez l'entreprise, d'après l'offre.]] Je suis très motivé à l'idée de [[deux ou trois responsabilités du poste]]. Mon autonomie, ma rigueur et ma capacité à apprendre rapidement me permettront de contribuer efficacement à votre équipe dès mon arrivée.

\vspace{10pt}

Je vous remercie de l'attention que vous porterez à ma candidature et reste à votre entière disposition pour un entretien.

\vspace{20pt}

Je vous prie d'agréer, Madame, Monsieur, l'expression de mes salutations distinguées.

\vspace{30pt}

\noindent [[MOI: Prénom Nom]]

\end{document}
